% Source: https://artofproblemsolving.com/community/c3h2516565_please_help_me_with_this_problem
% Stable identity: aops-topic-2516565-q01
Count unordered hands: there are $\dbinom{25}{4}=12650$ equally likely hands.

For a favourable hand, choose the $2$ numbers ($\dbinom52=10$ ways) and the $3$ colours ($\dbinom53=10$ ways). The hand then consists of $4$ of the $6$ cards in the $2\times3$ grid (numbers $\times$ colours), using all $3$ colours and both numbers. Four cards in three colour columns means one column contains both of its cards and the other two columns one card each. There are $3$ choices for the full column and $2\cdot2$ for the other two, i.e. $12$ ways, and both numbers then appear automatically (the full column has both). So there are $10\cdot10\cdot12=1200$ favourable hands, and
\[P=\dfrac{1200}{12650}=\boxed{\dfrac{24}{253}}.\]
(Counting ordered draws instead gives $25\cdot24\cdot23\cdot22=303600$ outcomes, and then each favourable hand must also be counted in all $4!=24$ orders: $\dfrac{1200\cdot24}{303600}=\dfrac{24}{253}$. Mixing an ordered total with unordered or partial counts, as in ``$25$, then $16$, then $4$, then $3$ choices'', is what goes wrong.)
